\chapter{偏振}
\problem{6-1}{自然光投射到互相重叠的两个偏振片上。如果透射光的强度为（1）透射光束最大强度的 $1/3$；（2）入射光束强度的 $1/3$，则这两个偏振片的透振方向之间夹角是多少？假定偏振片是理想的，它把自然光的强度严格减小一半。}
\answerarea
\problem{6-2}{一束自然光入射到一偏振片组上，这组由四块组成，每片的透振方向相对于前面一片沿顺时针方向转过 $30^\circ$ 角。试问入射光中有多大一部分透过了这组偏振片？}
\answerarea
\clearpage
\problem{6-3}{将一偏振片沿 $45^\circ$ 角插入一对正交偏振器之间，自然光经过它们强度减为原来的百分之几？}
\answerarea
\problem{6-4}{计算从空气到水面的布儒斯特角（水的折射率 $n=4/3$）。}
\answerarea
\clearpage
\problem{6-5}{一束光由水射在玻璃上，当入射角为 $50.82^\circ$ 时，反射光完全偏振。求玻璃的折射率（已知水的折射率为 $4/3$）。}
\answerarea
\problem{6-6}{计算：（1）由空气到玻璃（$n=1.560$）的全偏振角；（2）由此玻璃到空气的全偏振角；（3）在全偏振时由空气到此玻璃的折射光的偏振度；（4）在全偏振时由此玻璃到空气的折射光的偏振度；（5）自然光从空气以布儒斯特角入射到平行平面玻璃板以后，最终透射光的偏振度。}
\answerarea
\clearpage
\problem{6-7}{求自然光透过八块 $n=1.560$ 的平行玻璃板组成的玻片堆后的偏振度（忽略玻璃对光的吸收）。}
\answerarea
\problemfig[0.38]{6-8}{已知自然光射于某平行平面玻璃板上时，反射光的能流为入射光的 $0.10$ 倍（见本题图），取入射能流为一个单位，设玻璃的折射率为 $1.50$，求图中标出的光束2、3、4的能流（略去玻璃对光的吸收）。}{6_8}
\answerarea
\clearpage
\problem{6-9}{线偏振光的振动面和入射面之间的夹角称为振动的方位角。设入射线偏振光的方位角为 $\alpha$，入射角为 $i_1$，求折射光和反射光的方位角 $\alpha_2$ 和 $\alpha'_1$（已知两介质的折射率为 $n_1$ 和 $n_2$）。}
\answerarea
\problem{6-10}{线偏振光以布儒斯特角从空气入射到玻璃（$n=1.560$）的表面上，其振动方位角为 $20^\circ$，求反射光和折射光的方位角。}
\answerarea
\clearpage
\problem{6-11}{设入射光、反射光、折射光的总能流分别为 $W_1$、$W'_1$、$W_2$，则总能流反射率 $\mathcal R$ 和总能流透射率 $\mathcal T$ 定义为
\[\mathcal R=\frac{W'_1}{W_1},\qquad\mathcal T=\frac{W_2}{W_1}.\]
\begin{enumerate}\item 当入射光为线偏振光，方位角（见习题6-9）为 $\alpha$ 时，试证明
\[\mathcal R=\mathcal R_p\cos^2\alpha+\mathcal R_s\sin^2\alpha,\quad\mathcal T=\mathcal T_p\cos^2\alpha+\mathcal T_s\sin^2\alpha.\]
\item 证明 $\mathcal R+\mathcal T=1$。\item 设入射光是自然光，求 $\mathcal R$、$\mathcal T$ 与 $\mathcal R_p$、$\mathcal R_s$ 和 $\mathcal T_p$、$\mathcal T_s$ 的关系。\item 设入射光是圆偏振光，求 $\mathcal R$、$\mathcal T$ 与 $\mathcal R_p$、$\mathcal R_s$ 和 $\mathcal T_p$、$\mathcal T_s$ 的关系。\end{enumerate}}
\answerarea
\clearpage
\problem{6-12}{光从空气到玻璃（$n=1.50$）以布儒斯特角入射，试计算：（1）能流反射率 $\mathcal R_p$ 和 $\mathcal R_s$ 的值；（2）能流透射率 $\mathcal T_p$ 和 $\mathcal T_s$ 的值。}
\answerarea
\problem{6-13}{线偏振光从空气到玻璃（$n=1.5$）以 $45^\circ$ 角入射，方位角为 $60^\circ$，试计算：（1）总能流反射率 $\mathcal R$ 和总能流透射率 $\mathcal T$；（2）改为自然光入射，$\mathcal R$ 和 $\mathcal T$ 为多少？}
\answerarea
\clearpage
\problemfig[0.34]{6-14}{本题图所示为一支半导体砷化镓发光管。管芯 $AB$ 为发光区，其直径 $d\sim3\,\mathrm{mm}$。为了避免全反射，发光管上部研磨成半球形，以使内部发射光能够以最大的透射率向外输出。如果要使发光区边缘两点 $A$ 和 $B$ 发射的光不至于全反射，半球的半径至少应取多少？已知砷化镓的折射率为 $3.4$（对发射的 $0.9\,\mu\mathrm m$ 波长）。}{6_14}
\answerarea
\problemfig[0.34]{6-15}{接上题，为了减少光在砷化镓--空气界面的反射，工艺上常在砷化镓表面镀一层氧化硅增透膜。如本题图所示，氧化硅的折射率为 $1.7$。现在单纯从几何光学角度提出一个问题：加膜后入射角为多少才不致于在空气表面发生全反射？试与不加膜时相比（设膜很薄，可按平面板计算）。}{6_15}
\answerarea
\clearpage
\problem{6-16}{从光密介质到光疏介质，当 $\sin i_1>n_2/n_1$ 时就发生全反射。作为一种处理方法，我们仍可在形式上维持折射定律 $n_1\sin i_1=n_2\sin i_2$，这时 $\sin i_2>1$，可认为 $i_2$ 是个虚折射角，它的余弦也为虚数：
\[\cos i_2=\sqrt{1-\sin^2 i_2}=\mathrm i\sqrt{\left(\frac{n_1}{n_2}\right)^2\sin^2 i_1-1}.\]
试利用菲涅耳公式证明此时 $|r_p|=1$、$|r_s|=1$，从而 $\mathcal R_p=\mathcal R_s=1$。
\par【注：本来这里根号前有正、负两种可能。为了保证（6.55）式中的 $\kappa>0$，从而使（6.56）式中衰减因子的指数项是负的，根号应取正的。】}
\answerarea
\problem{6-17}{推导全反射时的相移公式（6.35）。
\begin{equation*}
\left\{\begin{aligned}
\delta_p&=2\arctan\!\left[
\frac{n_1}{n_2}\frac{\sqrt{\left(\dfrac{n_1}{n_2}\right)^2\sin^2 i_1-1}}{\cos i_1}
\right],\\[7pt]
\delta_s&=2\arctan\!\left[
\frac{n_2}{n_1}\frac{\sqrt{\left(\dfrac{n_1}{n_2}\right)^2\sin^2 i_1-1}}{\cos i_1}
\right].
\end{aligned}\right.
\tag{6.35}
\end{equation*}
适用条件为 $n_1>n_2$，且 $i_1>i_c=\arcsin(n_2/n_1)$。
\editorial{原题引用（6.34），该式实际为布儒斯特角公式；依据补充的原书页面，本题应引用全反射相移公式（6.35），已据此更正。}}
\answerarea
\clearpage
\problem{6-18}{（1）计算 $n_1=1.51$、$n_2=1.0$、入射角为 $54^\circ37'$ 时全反射光的相移 $\delta_p$ 和 $\delta_s$；（2）如果入射光是线偏振的，全反射光中 p 振动和 s 振动的相位差为多少？说明由两者合成的是椭圆偏振光。}
\answerarea
\problem{6-19}{在上题中用的光源是氦氖激光，求隐失波的穿透深度。}
\answerarea
\clearpage
\problem{6-20}{一束线偏振的钠黄光垂直射入一块方解石晶体，振动方向与晶体的主平面成 $20^\circ$ 角。试计算 o、e 两光束折射光的相对振幅。}
\answerarea
\problemfig[0.39]{6-21}{两大小相同的冰洲石晶体 $A$、$B$ 前后排列，强度为 $I_0$ 的自然光垂直于晶体 $A$ 的表面入射后相继通过 $A$、$B$（见本题图）。$A$、$B$ 的主截面之间的夹角为 $\alpha$（图中 $\alpha\ne0$）。求 $\alpha=0^\circ$、$45^\circ$、$90^\circ$、$180^\circ$ 时由 $B$ 射出光线的数目和每个的强度（忽略反射、吸收等损失）。}{6_21}
\answerarea
\clearpage
\problem{6-22}{一水晶平板厚 $0.850\,\mathrm{mm}$，光轴与表面平行，用水银灯的绿光（$546.1\,\mathrm{nm}$）垂直照射，求：（1）o、e 两光束在晶体中的光程；（2）二者的相位差（用度表示）。\missing{水晶对该波长的主折射率数据表未提供。}}
\answerarea
\problem{6-23}{一束钠黄光以 $50^\circ$ 的入射角射到冰洲石平板上。设光轴与板表面平行，并垂直于入射面，求晶体中 o 光和 e 光的夹角。}
\answerarea
\clearpage
\problem{6-24}{一束钠黄光掠入射到冰的晶体平板上，光轴与入射面垂直，平板厚度为 $4.20\,\mathrm{mm}$。求 o 光和 e 光射到平板对面上两点的间隔。已知对于钠黄光冰的 $n_o=1.3090$，$n_e=1.3104$。}
\answerarea
\problem{6-25}{用 ADP（磷酸二氢铵，$\mathrm{NH_4H_2PO_4}$）晶体制成 $50^\circ$ 顶角的棱镜，光轴与折射棱平行，$n_o=1.5246$，$n_e=1.4792$。试求：（1）o 光和 e 光的最小偏向角；（2）二者之差。}
\answerarea
\clearpage
\problem{6-26}{设一水晶棱镜的顶角为 $60^\circ$，光轴平行于折射棱。钠黄光以最小偏向角的方向在棱镜中折射，用焦距为 $1\,\mathrm m$ 的透镜聚焦，o 光和 e 光两谱线的间隔为多少？\missing{水晶对钠黄光的主折射率数据表未提供。}}
\answerarea
\problem{6-27}{求冰洲石晶体中光射线和波法线间的最大夹角。}
\answerarea
\clearpage
\problem{6-28}{图6-31b中，设入射光是钠黄光，晶体为方解石，光轴与晶体表面成 $30^\circ$ 角，入射角为 $45^\circ$。（1）o 光和 e 光的方向；（2）e 光的折射率。
\bookfig[0.7]{book_6_31b}{图6-31(b)\quad 用惠更斯作图法求单轴晶体中的折射线}
\par{\small 圆形与椭圆形次波面分别对应 o 光和 e 光，图中几何形状为示意。}
\missing{方解石对钠黄光的主折射率数据；图6-31b已补入。}}
\answerarea
\problem{6-29}{当图6-46中沃拉斯顿棱镜的顶角 $\alpha=15^\circ$ 时，两出射光线间的夹角为多少？
\bookfig[0.38]{book_6_46}{图6-46\quad 沃拉斯顿棱镜}
\par{\small 两块冰洲石棱镜的光轴方向互相垂直，图中以 $\updownarrow$、$\odot$ 示意；出射方向分别标为1和2。}}
\answerarea
\clearpage
\problem{6-30}{设图6-47b所示的尼科耳棱镜中 $\angle CA''C''$ 为直角，光线 $SM$ 平行于 $A''A$，计算 $\angle S_0MS$。
\bookfig[0.94]{book_6_47b}{图6-47(b)\quad 尼科耳棱镜的主截面}}
\answerarea
\problem{6-31}{用方解石和石英薄板作对钠黄光的 $\lambda/4$ 片，它们的最小厚度各为多少？}
\answerarea
\clearpage
\problem{6-32}{两尼科耳棱镜主截面的夹角由 $30^\circ$ 变到 $45^\circ$，透射光的强度如何变化？设入射自然光的强度为 $I_0$。}
\answerarea
\problem{6-33}{单色线偏振光垂直射入方解石晶体，其振动方向与主截面成 $30^\circ$ 角，两折射光线经过置于方解石后的尼科耳棱镜，其主截面与原入射光的振动方向成 $50^\circ$ 角。求两条光线的相对强度。}
\answerarea
\clearpage
\problem{6-34}{经尼科耳棱镜观察部分偏振光，当尼科耳棱镜由对应于极大强度的位置转过 $60^\circ$ 时，光强减为一半。求光束的偏振度。}
\answerarea
\problem{6-35}{用一块 $\lambda/4$ 片和一块偏振片鉴定一束椭圆偏振光。达到消光位置时，$\lambda/4$ 片的光轴与偏振片透振方向相差 $22^\circ$，求椭圆长短轴之比。}
\answerarea
\clearpage
\problem{6-36}{强度为 $I_0$ 的右旋圆偏振光垂直通过 $\lambda/4$ 片（此 $\lambda/4$ 片由方解石做成，o 光和 e 光在晶片中的光程差刚好是 $\lambda/4$），然后再经过一块主截面相对于 $\lambda/4$ 片光轴向右旋 $15^\circ$ 的尼科耳棱镜。求最后出射的光强（忽略反射、吸收等损失）。}
\answerarea
\problem{6-37}{平行于光轴切割一块方解石晶片，放置在主截面成 $35^\circ$ 角的一对尼科耳棱镜之间，晶片的光轴平分此角。求：（1）从方解石晶片射出的 o 光和 e 光的振幅和光强；（2）由第二个尼科耳棱镜射出的 o 光和 e 光的振幅和光强。设入射自然光的光强为 $I_0=A^2$，反射和吸收等损失可忽略。}
\answerarea
\clearpage
\problem{6-38}{强度为 $I_0$ 的单色平行光通过正交尼科耳棱镜。现在两尼科耳棱镜之间插入一 $\lambda/4$ 片，其主截面与第一尼科耳棱镜的主截面成 $60^\circ$ 角。求出射光的强度（忽略反射、吸收等损失）。}
\answerarea
\problem{6-39}{一块 $0.025\,\mathrm{mm}$ 厚的方解石晶片，表面平行于光轴，放在正交尼科耳棱镜之间，晶片的主截面与它们成 $45^\circ$ 角。试问：（1）在可见光范围内哪些波长的光不能通过？（2）如果将第二个尼科耳棱镜的主截面转到与第一个平行，哪些波长的光不能通过？}
\answerarea
\clearpage
\problem{6-40}{劈形水晶棱镜顶角为 $0.5^\circ$，棱边与光轴平行，置于正交尼科耳棱镜之间，使其主截面与两尼科耳棱镜的主截面都成 $45^\circ$ 角。以水银的 $404.7\,\mathrm{nm}$ 紫色平行光正入射。
\begin{enumerate}\item 通过第二尼科耳棱镜看到的干涉图样如何？\item 相邻暗纹的间隔 $d$ 等于多少？\item 若将第二尼科耳棱镜的主截面转 $90^\circ$，干涉图样有何变化？\item 维持两尼科耳棱镜正交，但让晶片的主截面转 $45^\circ$，使之与第二尼科耳棱镜的主截面垂直，干涉图样有何变化？\end{enumerate}}
\answerarea
\clearpage
\problem{6-41}{将比比涅补偿器（见思考题6-38）放在正交偏振片之间，光轴与它们的透振方向成 $45^\circ$ 角，你将看到什么现象？若劈角 $\alpha=2.75^\circ$，用平行的钠黄光照明，求干涉条纹的间隔。转动补偿器的光轴，对干涉条纹有什么影响？
\bookfig[0.65]{review_6_38}{思考题6-38图\quad 比比涅补偿器}
\par{\small 两块劈形棱镜的光轴相互垂直，分别以点状和横线纹理示意；$\alpha$ 为劈角。}}
\answerarea
\clearpage
\problem{6-42}{以线偏振光照在比比涅补偿器上（见思考题6-38），通过偏振片观察时在中央两劈形棱镜厚度 $d_1=d_2$ 处有一暗线，与中央暗线距离 $a$ 处又有一暗线。今以一同样波长的椭圆偏振光照在此比比涅补偿器上，发现暗线移至离中央 $b$ 处。
\begin{enumerate}\item 求椭圆偏振光在补偿器晶体中分解成的两个振动分量的初始相位差与 $a$、$b$ 的关系；\item 如果椭圆的长短轴正好分别与两劈形棱镜的光轴平行，试证此时 $b=a/4$；\item 设已知偏振片的透振方向与补偿器中一劈的光轴夹角为 $\theta$，找出 $\theta$ 与（2）问中椭圆长短轴比值的关系。\end{enumerate}
\bookfig[0.6]{review_6_38}{思考题6-38图\quad 比比涅补偿器（同6-41）}}
\answerarea
\clearpage
\problem{6-43}{已知水晶对钠黄光的旋光率 $\alpha=21.75^\circ/\mathrm{mm}$，求左、右旋圆偏振光折射率之差 $\Delta n$。}
\answerarea
\problem{6-44}{在两尼科耳棱镜之间插一块石英旋光晶片，以消除对眼睛最敏感的黄绿色光（$\lambda=550.0\,\mathrm{nm}$）。设对此波长的旋光率为 $24^\circ/\mathrm{mm}$，求下列情形下晶片的厚度：（1）两尼科耳棱镜主截面正交；（2）两尼科耳棱镜主截面平行。}
\answerarea
\clearpage
\problem{6-45}{一石英棒长 $5.639\,\mathrm{cm}$，端面垂直于光轴，置于正交偏振器间，沿光轴方向输入白光，用光谱仪观察透射光。
\begin{enumerate}\item 用一大张坐标纸，画出可见光范围（$400.0$--$760.0\,\mathrm{nm}$）振动面的旋转角与波长的曲线，旋光率数据可参照表6-5；\item 从这曲线看，哪些波长的光在光谱仪中消失？\item 在这些丢失的波长中，振动面的最大和最小旋转角各是多少？\end{enumerate}\missing{原书表6-5。}}
\answerarea
\clearpage
\problem{6-46}{一块表面垂直光轴的水晶片恰好抵消 $10\,\mathrm{cm}$ 长浓度20\%的麦芽糖溶液对钠光振动面所引起的旋转。对此波长水晶的旋光率 $\alpha=21.75^\circ/\mathrm{mm}$，麦芽糖的比旋光率 $[\alpha]=144^\circ/[\mathrm{dm}\cdot(\mathrm{g/cm^3})]$。求此水晶片的厚度。}
\answerarea
\problem{6-47}{ $15\,\mathrm{cm}$ 长的左旋葡萄糖溶液使钠黄光的振动面转了 $25.6^\circ$，已知 $[\alpha]=-51.4^\circ/[\mathrm{dm}\cdot(\mathrm{g/cm^3})]$，求溶液浓度。}
\answerarea
\clearpage
\problem{6-48}{将 $14.50\,\mathrm g$ 的蔗糖溶于水，得到 $60\,\mathrm{cm^3}$ 的溶液。在 $15\,\mathrm{cm}$ 的量糖计中测得钠光振动面旋转角为向右 $16.8^\circ$。已知 $[\alpha]=66.5^\circ/[\mathrm{dm}\cdot(\mathrm{g/cm^3})]$，这蔗糖样品中有多少比例的非旋光性杂质？}
\answerarea
\problem{6-49}{钠光以最小偏向角条件射入顶角为 $60^\circ$ 的石英晶体棱镜中，棱镜中光轴与底平行。求射出的左、右旋偏振光之间的夹角（所需数据在本章给出的表格中去找）。\missing{原书本章折射率、旋光率数据表。}}
\answerarea
\clearpage
